% ==========================================================================
% GENERATED by scripts/split_paper.py. Do not edit this file.
%
% Source: imported_papers/Koopman-JEPA-TMLR-rewrite/main.tex
% Part 7: Limitations
%
% Edits belong upstream, in the published paper. To change what the
% chapter includes, change imports/bounded.toml and re-run
% the splitter.
% ==========================================================================
\section{Limitations}
\label{sec:limitations}
% [draft 8 p1] DONE. Two changes to the draft. "We think this reading
% strengthens the case" becomes a statement of what the evidence shows, per the
% guide. And the draft named which control wins where, "mean-state on pendulum,
% persistence on cartpole", which is wrong on the seeded numbers: mean state
% wins three of the four tabulated cells and persistence the fourth. The
% sentence no longer needs the split, so it does not make the claim.
\paragraph{An Alternative Reading of the Gap.}
We read the disagreement of Section~\ref{sec:inversion} as a measurement artifact, and the opposite reading deserves stating in its strongest form. Every fixed readout of a latent state is fitted at some scale, whether it is a probe, a decoder, a value head or a controller, and a latent that leaves that scale is unusable to all of them however well it points. On that view the probe is not misreporting anything. It records a real failure that a downstream consumer would also suffer, and retrieval is the lenient metric, because cosine matching discards exactly the information such a consumer would need.

That reading strengthens the case for reporting the controls rather than weakening it. It is a claim about the readout and not about the dynamics, so it leaves the trivial predictors where they were. On the probe, the very metric this reading defends, a trivial predictor is the most accurate entry in both environments at both tabulated horizons (Table~\ref{tab:combined}). Accepting the alternative reading therefore relocates the failure rather than repairing it, from a metric that misreads the models to a metric on which the models lose to copying their own input. And if latent scale is what matters downstream, then controlling it is the objective, which Section~\ref{sec:spectrum} does at no training cost and slightly better than the architecture does.

% [orig, Empirical Scope] DONE, verbatim, with two sentences added on the
% behaviour policy. The draft's version of that addition said the random policy
% was used "rather than the intended mixture with an expert", which describes
% how the data came about rather than what it means for the claim. What matters
% to a reader is that the policy is random and that this is what narrows the
% pendulum distribution, so the sentences say that and nothing else.
% CONSUMED: the approved "Probe Choice" limitation paragraph is not ported
% here. Its content, that a two-layer probe absorbs coordinate changes a linear
% probe would penalize, belongs with the definition of the probe and lands in
% Section 4.3 at step 3.3.
\paragraph{Empirical Scope.}
Our environments are chosen for the falsifiability of the central spectral claim, not for visual complexity. The classical control benchmarks give us the ground-truth transparency to validate the bound of Equation~\ref{eq:linear_bound} against known physics. Scaling to richer environments (e.g., humanoid or wheeled-legged robots) raises distinct challenges. In contact-rich physics, the continuous dynamics are punctuated by discrete impacts, and a single global linear operator may struggle to capture these piecewise-linear regime shifts. A potential solution is to scale the latent dimension beyond $D{=}256$ and adopt a switching-Koopman architecture, where a learned gating mechanism selects from a library of specialized operators $\{\K_1, \dots, \K_M\}$ depending on the active interaction regime. Within these two environments the episodes are collected under a uniformly random policy, which is what narrows the pendulum state distribution at long horizons (Section~\ref{sec:degeneracy}). A mixture with an expert policy would broaden it, and would raise the level a trivial predictor has to clear.
% [draft 8 p3] DONE, trimmed. The draft's closing sentence said retrieval was
% computed only for the clean, full-data cell, so the later sections stayed
% probe-only. That caveat is gone because retrieval now exists in all 560
% cells, which is what Appendix E reports.
\paragraph{What Retrieval Cannot Do.}
Retrieval has biases of its own. It cannot resolve a prediction more finely than the spacing of the frames in its bank, so it compresses differences between models that are all accurate, and it flatters any model whose errors stay on the data manifold. It is a differently biased readout rather than a correct one, and the argument of Section~\ref{sec:inversion} rests on the disagreement between the two, not on either being authoritative.

% [new] DONE.
\paragraph{One Instance of a General Mechanism.}
The confound needs three ingredients: a prediction loss that is invariant to latent scale, a readout fitted at a fixed scale, and an autoregressive rollout that leaves the scale free to drift. All three are ordinary in the JEPA world-model line, and Section~\ref{sec:loss} identifies the loss as the lever, since changing that alone removes the gap. We have measured no pipeline but our own. Whether the same gap opens elsewhere is a question this design can pose and not answer, and the cheapest way to find out is the pairing of readouts in Section~\ref{sec:checklist}, which needs one extra evaluation pass and no retraining.

% =============================================================================
